\begin{abstractd}
Este reporte resuelve los cuatro ejercicios de la Actividad ponderada mediante integración por partes. En cada caso se conserva la elección indicada de $u$ y $dv$, se calculan $du$ y $v$, se aplica la fórmula $\int u\,dv=uv-\int v\,du$ y se verifica el resultado por diferenciación. El procedimiento muestra cómo integrar productos de funciones algebraicas con logaritmos, exponenciales y funciones trigonométricas.
\end{abstractd}

\templateIndex

\section{Introducción}
La integración por partes deriva de la regla del producto. Si
\[
 \frac{d}{dx}(uv)=u\frac{dv}{dx}+v\frac{du}{dx},
\]
al integrar y despejar se obtiene
\[
 \boxed{\int u\,dv=uv-\int v\,du}.
\]
El método es apropiado cuando el integrando es un producto y una de las funciones se simplifica al derivarse, mientras la otra puede integrarse directamente. En esta actividad las elecciones de $u$ y $dv$ ya están indicadas en la consigna; esto permite concentrarse en el cálculo correcto de $du$, $v$ y la integral restante \citep{stewart2016calculus,openstax2016calculus2}.

\section{Fórmulas y procedimiento resumido}
\subsection{Fórmulas auxiliares}
\[
 \frac{d}{dx}\ln x=\frac1x\quad(x>0),\qquad
 \int x^n\,dx=\frac{x^{n+1}}{n+1}+C\quad(n\neq-1),
\]
\[
 \int e^x\,dx=e^x+C,
 \qquad \int\sen(ax)\,dx=-\frac{\cos(ax)}{a}+C,
 \qquad \int\cos(ax)\,dx=\frac{\sen(ax)}{a}+C.
\]

\subsection{Secuencia de trabajo}
\begin{enumerate}
  \item Identificar $u$ y $dv$ según la consigna.
  \item Derivar $u$ para obtener $du$ e integrar $dv$ para obtener $v$.
  \item Sustituir en $\int u\,dv=uv-\int v\,du$.
  \item Resolver la integral restante, simplificar y añadir $C$.
  \item Derivar la respuesta y comprobar que se recupera el integrando.
\end{enumerate}

\section{Resolución de ejercicios}

\begin{examplebox}{Ejercicio 7: $\int x^3\ln x\,dx$}
\textbf{Elección indicada:}
\[
 u=\ln x,\qquad dv=x^3\,dx.
\]
Derivamos $u$ e integramos $dv$:
\[
 du=\frac{d}{dx}(\ln x)\,dx=\frac{1}{x}\,dx,
 \qquad
 v=\int x^3dx=\frac{x^{3+1}}{3+1}=\frac{x^4}{4}.
\]
Sustituimos $u$, $v$ y $du$ en $uv-\int v\,du$:
\begin{align*}
 \int x^3\ln x\,dx
 &=(\ln x)\left(\frac{x^4}{4}\right)
   -\int\left(\frac{x^4}{4}\right)\left(\frac{1}{x}\,dx\right)\\
 &=\frac{x^4}{4}\ln x-\int\frac{x^4}{4x}\,dx\\
 &=\frac{x^4}{4}\ln x-\frac14\int x^3dx\\
 &=\frac{x^4}{4}\ln x-\frac14\left(\frac{x^4}{4}\right)+C\\
 &=\frac{x^4}{4}\ln x-\frac{x^4}{16}+C.
\end{align*}
\textbf{Resultado:}
\[
 \boxed{\int x^3\ln x\,dx=\frac{x^4}{4}\ln x-\frac{x^4}{16}+C},\qquad x>0.
\]
\textbf{Verificación:}
\begin{align*}
 \frac{d}{dx}\left(\frac{x^4}{4}\ln x-\frac{x^4}{16}\right)
 &=\frac{d}{dx}\left(\frac{x^4}{4}\ln x\right)
   -\frac{d}{dx}\left(\frac{x^4}{16}\right)\\
 &=\left(x^3\ln x+\frac{x^4}{4}\frac1x\right)-\frac{4x^3}{16}\\
 &=x^3\ln x+\frac{x^3}{4}-\frac{x^3}{4}=x^3\ln x.
\end{align*}
\end{examplebox}

\begin{examplebox}{Ejercicio 8: $\int(4x+7)e^x\,dx$}
\textbf{Elección indicada:}
\[
 u=4x+7,\qquad dv=e^x\,dx.
\]
Derivamos $u$ e integramos $dv$:
\[
 du=\frac{d}{dx}(4x+7)\,dx=4\,dx,
 \qquad v=\int e^x\,dx=e^x.
\]
Sustituimos en $uv-\int v\,du$:
\begin{align*}
 \int(4x+7)e^x\,dx
 &=(4x+7)e^x-\int e^x(4\,dx)\\
 &=(4x+7)e^x-4\int e^x\,dx\\
 &=(4x+7)e^x-4e^x+C\\
 &=\bigl[(4x+7)-4\bigr]e^x+C\\
 &=(4x+3)e^x+C.
\end{align*}
\textbf{Resultado:}
\[
 \boxed{\int(4x+7)e^x\,dx=(4x+3)e^x+C}.
\]
\textbf{Verificación:}
\begin{align*}
 \frac{d}{dx}\bigl((4x+3)e^x\bigr)
 &=(4x+3)'e^x+(4x+3)(e^x)'\\
 &=4e^x+(4x+3)e^x\\
 &=(4x+7)e^x.
\end{align*}
\end{examplebox}

\begin{examplebox}{Ejercicio 9: $\int x\sen(3x)\,dx$}
\textbf{Elección indicada:}
\[
 u=x,\qquad dv=\sen(3x)\,dx.
\]
Derivamos $u$. Para integrar $dv$, usamos $w=3x$, de modo que
$dw=3\,dx$ y $dx=dw/3$:
\[
 du=\frac{d}{dx}(x)\,dx=dx,
 \qquad
 v=\int\sen(3x)dx
 =\frac13\int\sen w\,dw
 =-\frac13\cos(3x).
\]
Sustituimos en $uv-\int v\,du$:
\begin{align*}
 \int x\sen(3x)\,dx
 &=x\left(-\frac13\cos(3x)\right)
   -\int\left(-\frac13\cos(3x)\right)dx\\
 &=-\frac{x}{3}\cos(3x)+\frac13\int\cos(3x)dx\\
 &=-\frac{x}{3}\cos(3x)
   +\frac13\left(\frac13\sen(3x)\right)+C\\
 &=-\frac{x}{3}\cos(3x)+\frac19\sen(3x)+C.
\end{align*}
\textbf{Resultado:}
\[
 \boxed{\int x\sen(3x)\,dx=-\frac{x}{3}\cos(3x)+\frac19\sen(3x)+C}.
\]
\textbf{Verificación:}
\begin{align*}
 \frac{d}{dx}\left(-\frac{x}{3}\cos(3x)+\frac19\sen(3x)\right)
 &=-\frac13\left[x'\cos(3x)+x\bigl(\cos(3x)\bigr)'\right]
   +\frac19\bigl(\sen(3x)\bigr)'\\
 &=-\frac13\left[\cos(3x)-3x\sen(3x)\right]
   +\frac19\left[3\cos(3x)\right]\\
 &=-\frac13\cos(3x)+x\sen(3x)+\frac13\cos(3x)\\
 &=x\sen(3x).
\end{align*}
\end{examplebox}

\begin{examplebox}{Ejercicio 10: $\int x\cos(4x)\,dx$}
\textbf{Elección indicada:}
\[
 u=x,\qquad dv=\cos(4x)\,dx.
\]
Derivamos $u$. Para integrar $dv$, hacemos $w=4x$; entonces
$dw=4\,dx$ y $dx=dw/4$:
\[
 du=\frac{d}{dx}(x)\,dx=dx,
 \qquad
 v=\int\cos(4x)dx
 =\frac14\int\cos w\,dw
 =\frac14\sen(4x).
\]
Sustituimos en $uv-\int v\,du$:
\begin{align*}
 \int x\cos(4x)\,dx
 &=x\left(\frac14\sen(4x)\right)
   -\int\left(\frac14\sen(4x)\right)dx\\
 &=\frac{x}{4}\sen(4x)-\frac14\int\sen(4x)dx\\
 &=\frac{x}{4}\sen(4x)
   -\frac14\left(-\frac14\cos(4x)\right)+C\\
 &=\frac{x}{4}\sen(4x)+\frac1{16}\cos(4x)+C.
\end{align*}
\textbf{Resultado:}
\[
 \boxed{\int x\cos(4x)\,dx=\frac{x}{4}\sen(4x)+\frac1{16}\cos(4x)+C}.
\]
\textbf{Verificación:}
\begin{align*}
 \frac{d}{dx}\left(\frac{x}{4}\sen(4x)+\frac1{16}\cos(4x)\right)
 &=\frac14\left[x'\sen(4x)+x\bigl(\sen(4x)\bigr)'\right]
   +\frac1{16}\bigl(\cos(4x)\bigr)'\\
 &=\frac14\left[\sen(4x)+4x\cos(4x)\right]
   +\frac1{16}\left[-4\sen(4x)\right]\\
 &=\frac14\sen(4x)+x\cos(4x)-\frac14\sen(4x)\\
 &=x\cos(4x).
\end{align*}
\end{examplebox}

\section{Comprobación global}
Las cuatro verificaciones producen exactamente los integrandos originales. Además, una comprobación simbólica independiente confirmó que la diferencia entre la derivada de cada respuesta y su integrando es cero. La constante $C$ aparece en todas las respuestas porque se trata de integrales indefinidas.

\section{Conclusiones}
La integración por partes transforma un producto difícil en otro más sencillo. En el ejercicio logarítmico, derivar $\ln x$ elimina el logaritmo; en el exponencial, derivar el factor lineal lo convierte en una constante; y en los ejercicios trigonométricos, derivar $x$ reduce el producto a una integral inmediata. Los signos y factores $1/3$ y $1/4$ son esenciales al integrar funciones trigonométricas compuestas. La verificación por derivación confirma los cuatro resultados y constituye el control más directo del procedimiento.

\nocite{thomas2018calculus,openstax2016calculus1}
\bibliography{UANL/ingeniero-agronomo/calculo-integral/calculo-integral}
